% !PRINTED 519
\weberpage{519}
%<<CONT
so muss jedes $\varphi(\alpha_i, y_\alpha)$ gleich Null sein, d.\,h. $\varphi(x, y_\alpha)$ muss
durch $f(x, y_\alpha)$ theilbar sein. Es folgt aber hieraus nicht, dass
$f(x, y_\alpha) = f_\alpha(x, \psi)$ auch im K\"orper $\Omega(\psi)$ irreducibel ist.

Dies sind dieselben S\"atze, wie im \S.\,151. Es ist dort noch
gezeigt, dass man f\"ur $y_\alpha$ eine Function von $\alpha$ allein nehmen kann,
worauf wir nicht nochmals eingehen wollen.

Von Interesse ist es aber, den Zusammenhang der \gesperrt{Galois}'schen
Gruppen der einzelnen Factoren $f(x, y_\alpha), f(x, y_\beta), \ldots f(x, y_\sigma)$
zu verfolgen. Wir haben schon oben gesehen, dass man die
Gruppe des Factors $f(x, y_\alpha)$ im K\"orper $\Omega(\psi)$ erh\"alt, wenn man
alle Permutationen der $\alpha$ aufsucht, die durch $Q$ hervorgerufen
werden. Wir wollen diese Gruppe mit $P_\alpha$ bezeichnen, und ebenso
f\"ur die anderen Factoren die Zeichen $P_\beta, \ldots P_\sigma$ gebrauchen.

Wir greifen aus $P$ irgend $s - 1$ Permutationen $\pi_\beta \ldots \pi_\sigma$
von der Art heraus, dass durch $\pi_\beta$ die Reihe $A$ in $B$ etc.
durch $\pi_\sigma$ die Reihe $A$ in $S$ \"ubergeht, und wir wollen nun der
Einfachheit halber, was offenbar noch vollkommen freisteht, die
Bezeichnung der $\beta, \ldots \sigma$ so gew\"ahlt annehmen, dass $\pi_\beta, \ldots \pi_\sigma$
folgende Vertauschungen hervorrufen:
\[
\begin{pmatrix} \alpha_0, & \alpha_1 & \ldots & \alpha_{r-1} \\ \beta_0, & \beta_1 & \ldots & \beta_{r-1} \end{pmatrix}, \ldots
\begin{pmatrix} \alpha_0, & \alpha_1 & \ldots & \alpha_{r-1} \\ \sigma_0, & \sigma_1 & \ldots & \sigma_{r-1} \end{pmatrix}.
\]
Es sei nun $\varkappa$ irgend ein Element von $Q$, das unter den $\alpha$ die
Vertauschung
\[
(\alpha) = \begin{pmatrix} \alpha_0, & \alpha_1 & \ldots & \alpha_{r-1} \\ \alpha_{a_0}, & \alpha_{a_1} & \ldots & \alpha_{a_{r-1}} \end{pmatrix}
\]
bewirkt, dann ist
\[
\varkappa_\beta = \pi_\beta^{-1} \varkappa \pi_\beta
\]
ebenfalls in $Q$ enthalten, und es wird durch $\varkappa_\beta$ unter den $\beta$ die
Vertauschung
\[
(\beta) = \begin{pmatrix} \beta_0, & \beta_1 & \ldots & \beta_{r-1} \\ \beta_{a_0}, & \beta_{a_1} & \ldots & \beta_{a_{r-1}} \end{pmatrix},
\]
also unter den Indices der $\beta$ genau dieselbe Permutation hervorgerufen,
wie durch $\varkappa$ unter den Indices von $\alpha$. Nehmen wir
$(\beta)$ und ein zugeh\"origes $\varkappa_\beta$ als gegeben an, so ergiebt $\pi_\beta \varkappa_\beta \pi_\beta^{-1}$
wieder $\varkappa$, so dass also die durch $Q$ hervorgerufenen Permutationen
 der $\beta$ genau \"ubereinstimmen mit denen, die unter den
$\alpha$ hervorgerufen werden; und dasselbe gilt f\"ur die \"ubrigen
Reihen. Damit ist aber der Satz bewiesen.
